\honest{Entropy, and Short Codes}{Eliezer Yudkowsky}{2008}

A system with eight equally likely states has an entropy of 3 bits: you need three yes-or-no questions to find its state. A system Y whose four states have probabilities 1/2, 1/4, 1/8 and 1/8 has an entropy of 1.75 bits. I show what 1.75 questions means with the code 1, 01, 001, 000: half the time one question settles it, a quarter of the time two, and otherwise three. The general formula is the sum of $-p \log_2 p$ over the states. \nb{All the numbers check. The formula is Claude Shannon's, from 1948.}

A perfect code is not always possible, I say, but coding many copies at once comes arbitrarily close.

Why not give Y4 the shorter word 10? Because then ``1001'' could mean Y4 followed by Y2, or Y1 followed by Y3. ``The moral is that short words are a conserved resource.'' \nb{The exact form of this is the Kraft--McMillan inequality. The post's code already uses the whole budget, so no word can be shortened without lengthening another.}

In the limit, the length of a message ``corresponds exactly or almost exactly to its probability.'' This, I say, is the Minimum Description Length or Minimum Message Length form of Occam's razor. ``And so even the labels that we use for words are not quite arbitrary.'' \nb{A word's length in a language is set by how often the word is used; Minimum Message Length measures how probable a hypothesis is. Yudkowsky kept the two apart in ``Occam's Razor,'' where ``witch'' is short but names ``extraordinary assertions,'' and again in ``37 Ways That Words Can Be Wrong.''} \nb{The arbitrariness linguists assert, after Saussure, concerns sound and meaning, and is compatible with this point about length.}

The idea that ``You can X any way you like,'' I say, is ``a huge obstacle to learning how to X wisely.'' ``I can define a word any way I like'' obstructs carving reality at its joints, and even the sensible-sounding ``The labels we attach to words are arbitrary'' obstructs awareness of compactness. Tolkien heard how beautiful ``cellar door'' sounds; that, I say, is the kind of awareness it takes to use language like Tolkien.''

``Chair,'' in the terminology of cognitive psychology, is a basic-level category: people talk about chairs, rather than recliners or furniture. \nb{Rosch and colleagues found in 1976 that people naming pictured objects ``overwhelmingly used the basic level name.''} It is no coincidence, I say, that ``chair'' has fewer syllables than the other two: ``Frequent use goes along with short words; short words go along with frequent use.'' \nb{The general rule is Zipf's law of abbreviation, confirmed in almost a thousand languages. The route through basic-level categories is less direct: in Rosch's study the more general word was the more frequent one in nine of fifteen comparisons.}

I close with a joke attributed to Douglas Hofstadter about the words ``the'' and ``antidisestablishmentarianism.''
